\section{Data of the examples}\label{app:data}
This appendix lists the data of the examples of Section~\ref{sec:examples}: for $X_9$, $X_{19}$ and $X_{20}$ the rays of
$\Sigma'$, the node parallelograms, a $\Z$-basis of $K$ and the tropical $2$-cycles of Lemma~\ref{lem:x9belt} and
Table~\ref{tab:x19x20}; for $X_{88}$ and $X_{154}$ the vertices, the profile, the node parallelograms and a $\Z$-basis of
$K$. The accompanying computation bundle contains the triangulations $\mathcal T$, integral heights, fans $\Sigma'$,
polyhedral bases and cycle triangles for all five examples. Its checker recomputes the toric data, relation lattices,
boundary networks and coefficients, and the stated Euler characteristics and numbers of interior intersections with
$\Gamma$. The printed data below are compared with these records. The checker does not reconstruct the coupled base
from its two fans or establish every interior condition of \cite[Def.~7.2]{CBM}. For the belt cycles it checks that each
is a disc with one field that meets $\Gamma$ only along its boundary, and for the cycle of $X_9$ the conditions of
\cite[Def.~7.2]{CBM} are proved in Lemma~\ref{lem:x9belt}; for the other cycles, those with trivalent vertices and the
cycle of $X_{88}$ whose boundary is a closed path and which crosses $\Gamma$ at $18$ interior points, the interior
conditions rest on the computations described in Section~\ref{ssec:evidence}.

\emph{Conventions.} The rays of $\Sigma'$ are numbered $n_0,n_1,\dots$; the vertices of $\Delta$ come first. A node
parallelogram is written $q=(a,b,c,d)$ with $n_a+n_c=n_b+n_d$, in the labelling of Section~\ref{ssec:admissible}:
$\{a,c\}$ is the chosen diagonal, which passes through the centre at a pentagon or a hexagon, and
$\circv_q=e_b+e_d-e_a-e_c$. A \emph{boundary segment} of a tropical $2$-cycle $S$ is a maximal path in $\partial S$ whose
interior avoids the negative vertices; it is either closed or joins two negative vertices $t_i$. A boundary segment is
written, in one of its two directions, as the sequence of the lower supports of its arcs and of the nodes it passes. An
ordered pair $(i,j)$ stands for the arcs with lower support $\{n_i,n_j\}$ on which the boundary network $\partial(v[S])$,
read in the direction of listing, has coefficient $n_j-n_i$; so the boundary field is $v=n_j-n_i$ if the listing follows
the boundary orientation of $\partial S$ and $v=n_i-n_j$ otherwise, and $e_{n_j}-e_{n_i}$ is the corresponding canonical
lift (Section~\ref{ssec:tropical}). Consecutive arcs with the same pair, which meet at positive vertices or at edge
points, are listed once, and $q_k$ between two pairs marks the passage through the node $q_k$. A closed boundary segment
returns to its first pair after the last node listed. Reversing the direction of listing reverses the order and the
signs, so the difference of the lifts after and before $q_k$ does not depend on it; by Definition~\ref{def:nodecoef},
$\rho_{q_k}(S)$ is the coefficient of $\circv_{q_k}$ in this difference, and can therefore be read off from the lists
below. For the cycles with trivalent vertices, we also give the number of four-valent interior vertices of $S$, the
number of points at which $S$ crosses $\Gamma$ transversally (clause (b) of \cite[Def.~7.2]{CBM}), and the Euler
characteristic of $S$. Vectors in $\Z^Q$ are written as combinations of the nodes: $q_2-q_3$ stands for the vector with
coordinate $1$ at $q_2$, $-1$ at $q_3$ and $0$ elsewhere.

\begingroup\sloppy
\subsection{The example \texorpdfstring{$X_9$}{X9}}\label{app:data-x9}
\emph{Rays of $\Sigma'$.} The vertices of $\Delta_9$ are $n_{0}=(1,0,0,0)$, $n_{1}=(0,1,0,0)$, $n_{2}=(-1,-1,0,0)$, $n_{3}=(0,0,1,0)$, $n_{4}=(0,0,0,1)$, $n_{5}=(-4,-2,-1,0)$, $n_{6}=(-4,-2,0,-1)$, $n_{7}=(3,1,1,1)$, $n_{8}=(2,1,1,1)$.
 The centre of the $dP_7$ pentagon is $n_{10}=(-2,-1,0,0)$. The fan $\Sigma'$ has one further ray, $n_{9}=(-1,0,0,0)$, in the interior of a facet of $\Delta_9$.

\emph{Node parallelograms} $q=(a,b,c,d)$, given by the indices of $n_a,n_b,n_c,n_d$; an asterisk marks the two parallelograms of the pentagon: $q_{1}=(3,6,10,8)^{*}$, $q_{2}=(2,6,3,7)$, $q_{3}=(4,5,10,8)^{*}$, $q_{4}=(2,5,4,7)$.

\emph{A $\Z$-basis of $K$} (rank 1): $q_{1}-q_{2}-q_{3}+q_{4}$.

\emph{Generating cycles.}
The cycles are the following.
\begin{itemize}
\item $S_{1}$: a belt cycle with constant field $\pm(4,2,1,1)$, a disc; its boundary passes 5 positive vertices.
Boundary segments: $(3,6)$ $q_{2}$ $(7,2)$ $q_{4}$ $(4,5)$ $q_{3}$ $(8,10)$ $q_{1}$.
Node coefficients: $\rho(S_{1})=q_{1}-q_{2}-q_{3}+q_{4}$.
\end{itemize}
\emph{Index.} The node coefficient vector of $S_1$ is the basis vector of $K$.

\subsection{The example \texorpdfstring{$X_{19}$}{X19}}\label{app:data-x19}
\emph{Rays of $\Sigma'$.} The vertices of $\Delta_{19}$ are $n_{0}=(1,0,0,0)$, $n_{1}=(0,1,0,0)$, $n_{2}=(0,0,1,0)$, $n_{3}=(0,0,-1,0)$, $n_{4}=(0,-1,0,0)$, $n_{5}=(1,-1,-1,0)$, $n_{6}=(0,0,0,1)$, $n_{7}=(-1,1,1,-1)$, $n_{8}=(0,-1,0,1)$, $n_{9}=(0,0,-1,1)$, $n_{10}=(-1,1,0,-1)$, $n_{11}=(-1,0,1,-1)$, $n_{12}=(-1,0,1,0)$, $n_{13}=(-1,1,0,0)$, $n_{14}=(0,-1,-1,0)$, $n_{15}=(0,-1,-1,1)$, $n_{16}=(-1,0,0,-1)$, $n_{17}=(-2,1,1,-1)$, $n_{18}=(-1,0,0,1)$.
 The centre of the $dP_7$ pentagon is $n_{19}=(-1,0,0,0)$.

\emph{Node parallelograms} $q=(a,b,c,d)$, given by the indices of $n_a,n_b,n_c,n_d$; an asterisk marks the two parallelograms of the pentagon: $q_{1}=(16,17,19,14)^{*}$, $q_{2}=(1,7,17,13)$, $q_{3}=(2,7,17,12)$, $q_{4}=(2,6,18,12)$, $q_{5}=(15,18,19,14)^{*}$, $q_{6}=(4,8,12,11)$, $q_{7}=(1,6,18,13)$, $q_{8}=(4,11,16,14)$, $q_{9}=(3,10,16,14)$, $q_{10}=(3,9,13,10)$, $q_{11}=(4,8,15,14)$, $q_{12}=(3,9,15,14)$, $q_{13}=(0,2,4,5)$, $q_{14}=(7,10,16,11)$, $q_{15}=(6,8,15,9)$, $q_{16}=(0,1,3,5)$.

\emph{A $\Z$-basis of $K$} (rank 3): $q_{2}-q_{3}+q_{4}-q_{7}$, $q_{2}-q_{3}-q_{6}+q_{8}-q_{9}+q_{10}+q_{11}-q_{12}-q_{13}+q_{16}$, $q_{1}+q_{2}-q_{5}-q_{7}-q_{8}-q_{9}+q_{11}+q_{12}+q_{14}-q_{15}$.

\emph{Generating cycles.}
The boundary graphs below have the following negative vertices.
\begin{itemize}
\item $t_{1}$: lower support $\{n_{10},n_{16},n_{17}\}$, face label $[(0,0,0,1),(1,0,1,0)]$.
\item $t_{2}$: lower support $\{n_{7},n_{10},n_{17}\}$, face label $[(0,-1,0,0),(0,0,0,1)]$.
\item $t_{3}$: lower support $\{n_{8},n_{15},n_{18}\}$, face label $[(0,0,0,-1),(1,1,0,0)]$.
\item $t_{4}$: lower support $\{n_{6},n_{8},n_{18}\}$, face label $[(0,0,-1,-1),(0,0,0,-1)]$.
\item $t_{5}$: lower support $\{n_{2},n_{11},n_{12}\}$, face label $[(0,0,-1,0),(0,1,-1,0)]$.
\item $t_{6}$: lower support $\{n_{1},n_{9},n_{13}\}$, face label $[(0,-1,0,-1),(0,-1,1,0)]$.
\item $t_{7}$: lower support $\{n_{2},n_{4},n_{11}\}$, face label $[(-1,1,-1,1),(0,1,-1,0)]$.
\item $t_{8}$: lower support $\{n_{1},n_{3},n_{9}\}$, face label $[(-1,-1,1,0),(0,-1,1,0)]$.
\end{itemize}
The cycles are the following.
\begin{itemize}
\item $S_{1}$: a belt cycle with constant field $\pm(1,0,0,0)$, a disc; its boundary passes 5 positive vertices.
Boundary segments: $(17,7)$ $q_{3}$ $(12,2)$ $q_{4}$ $(18,6)$ $q_{7}$ $(13,1)$ $q_{2}$.
Node coefficients: $\rho(S_{1})=q_{2}-q_{3}+q_{4}-q_{7}$.
\item $S_{2}$: boundary graph a complete graph on four vertices; 20 positive vertices on $\partial S$, 6 four-valent interior vertices, 11 transverse crossings with $\Gamma$, $\chi(S)=5$.
Boundary segments: $t_{1}$ $(17,16)$ $q_{1}$ $(19,14)$ $q_{5}$ $(18,15)$ $t_{3}$; $t_{1}$ $(10,17)$ $t_{2}$; $t_{1}$ $(16,10)$ $q_{9}$ $(14,3)$ $q_{12}$ $(15,9)$ $q_{15}$ $(8,6)$ $t_{4}$; $t_{2}$ $(7,17)$ $q_{3}$ $(2,12)$ $q_{4}$ $(6,18)$ $t_{4}$; $t_{2}$ $(10,7)$ $q_{14}$ $(16,11)$ $q_{8}$ $(14,4)$ $q_{11}$ $(15,8)$ $t_{3}$; $t_{3}$ $(18,8)$ $t_{4}$.
Node coefficients: $\rho(S_{2})=q_{1}+q_{3}-q_{4}-q_{5}-q_{8}-q_{9}+q_{11}+q_{12}+q_{14}-q_{15}$.
\item $S_{3}$: boundary graph a complete graph on four vertices; 26 positive vertices on $\partial S$, 17 four-valent interior vertices, 39 transverse crossings with $\Gamma$, $\chi(S)=5$.
Boundary segments: $t_{5}$ $(12,11)$ $q_{6}$ $(8,4)$ $q_{11}$ $(15,14)$ $q_{12}$ $(9,3)$ $t_{8}$; $t_{5}$ $(11,2)$ $t_{7}$; $t_{5}$ $(2,12)$ $q_{3}$ $(7,17)$ $q_{2}$ $(1,13)$ $t_{6}$; $t_{6}$ $(9,13)$ $q_{10}$ $(3,10)$ $q_{9}$ $(14,16)$ $q_{8}$ $(4,11)$ $t_{7}$; $t_{6}$ $(1,9)$ $t_{8}$; $t_{7}$ $(4,2)$ $q_{13}$ $(5,0)$ $q_{16}$ $(3,1)$ $t_{8}$.
Node coefficients: $\rho(S_{3})=q_{2}-q_{3}-q_{6}+q_{8}-q_{9}+q_{10}+q_{11}-q_{12}-q_{13}+q_{16}$.
\end{itemize}
\emph{Index.} In the basis above, the node coefficient vectors of $S_1,\dots,S_{3}$ have coordinates $(1,0,0)$, $(-1,0,1)$, $(0,1,0)$; the elementary divisors of this matrix are $1,1,1$, so the vectors $\rho(S_i)$ generate $K$.

\subsection{The example \texorpdfstring{$X_{20}$}{X20}}\label{app:data-x20}
\emph{Rays of $\Sigma'$.} The vertices of $\Delta_{20}$ are $n_{0}=(1,0,0,0)$, $n_{1}=(0,1,0,0)$, $n_{2}=(0,0,1,0)$, $n_{3}=(0,0,0,1)$, $n_{4}=(1,-1,-1,-1)$, $n_{5}=(-1,1,1,1)$, $n_{6}=(0,0,0,-1)$, $n_{7}=(0,0,-1,0)$, $n_{8}=(0,-1,0,0)$, $n_{9}=(-1,1,0,1)$, $n_{10}=(0,0,-1,-1)$, $n_{11}=(-1,1,0,0)$, $n_{12}=(-1,0,1,1)$, $n_{13}=(0,-1,0,-1)$, $n_{14}=(-1,0,1,0)$, $n_{15}=(0,-1,-1,-1)$, $n_{16}=(-1,0,0,1)$, $n_{17}=(-1,0,-1,0)$, $n_{18}=(-1,-1,0,0)$, $n_{19}=(-1,-1,1,0)$.
 The centre of the $dP_7$ pentagon is $n_{20}=(-1,0,0,0)$.

\emph{Node parallelograms} $q=(a,b,c,d)$, given by the indices of $n_a,n_b,n_c,n_d$; an asterisk marks the two parallelograms of the pentagon: $q_{1}=(18,19,20,17)^{*}$, $q_{2}=(13,15,18,19)$, $q_{3}=(12,16,18,19)$, $q_{4}=(3,8,18,16)$, $q_{5}=(4,8,18,15)$, $q_{6}=(14,19,20,11)^{*}$, $q_{7}=(6,13,19,14)$, $q_{8}=(5,12,19,14)$, $q_{9}=(4,7,17,15)$, $q_{10}=(3,7,17,16)$, $q_{11}=(5,9,16,12)$, $q_{12}=(6,10,15,13)$, $q_{13}=(4,8,19,13)$, $q_{14}=(3,8,19,12)$, $q_{15}=(0,1,9,3)$, $q_{16}=(0,2,19,8)$, $q_{17}=(0,2,12,3)$, $q_{18}=(0,1,10,4)$, $q_{19}=(0,2,13,4)$.

\emph{A $\Z$-basis of $K$} (rank 6): $q_{14}-q_{16}+q_{17}$, $q_{13}-q_{16}+q_{19}$, $q_{3}-q_{4}+q_{14}$, $q_{2}-q_{5}+q_{13}$, $q_{4}-q_{5}+q_{9}-q_{10}$, $q_{4}-q_{5}+q_{7}-q_{8}+q_{11}-q_{12}+q_{15}-q_{18}$.

\emph{Generating cycles.}
The boundary graphs below have the following negative vertices.
\begin{itemize}
\item $t_{1}$: lower support $\{n_{3},n_{9},n_{16}\}$, face label $[(0,0,0,-1),(0,0,1,-1)]$.
\item $t_{2}$: lower support $\{n_{4},n_{10},n_{15}\}$, face label $[(0,0,0,1),(0,0,1,0)]$.
\end{itemize}
The cycles are the following.
\begin{itemize}
\item $S_{1}$: a belt cycle with constant field $\pm(1,0,-1,0)$, a disc; its boundary passes 3 positive vertices.
Boundary segments: $(8,19)$ $q_{16}$ $(0,2)$ $q_{17}$ $(3,12)$ $q_{14}$.
Node coefficients: $\rho(S_{1})=-q_{14}+q_{16}-q_{17}$.
\item $S_{2}$: a belt cycle with constant field $\pm(0,1,0,1)$, a disc; its boundary passes 5 positive vertices.
Boundary segments: $(18,16)$ $q_{4}$ $(8,3)$ $q_{14}$ $(19,12)$ $q_{3}$.
Node coefficients: $\rho(S_{2})=q_{3}-q_{4}+q_{14}$.
\item $S_{3}$: a belt cycle with constant field $\pm(1,0,0,0)$, a disc; its boundary passes 5 positive vertices.
Boundary segments: $(8,18)$ $q_{5}$ $(4,15)$ $q_{9}$ $(7,17)$ $q_{10}$ $(3,16)$ $q_{4}$.
Node coefficients: $\rho(S_{3})=-q_{4}+q_{5}-q_{9}+q_{10}$.
\item $S_{4}$: a belt cycle with constant field $\pm(1,0,-1,-1)$, a disc; its boundary passes 5 positive vertices.
Boundary segments: $(18,15)$ $q_{5}$ $(8,4)$ $q_{13}$ $(19,13)$ $q_{2}$.
Node coefficients: $\rho(S_{4})=q_{2}-q_{5}+q_{13}$.
\item $S_{5}$: a belt cycle with constant field $\pm(1,0,-1,0)$, a disc; its boundary passes 6 positive vertices.
Boundary segments: $(8,19)$ $q_{16}$ $(0,2)$ $q_{19}$ $(4,13)$ $q_{13}$.
Node coefficients: $\rho(S_{5})=-q_{13}+q_{16}-q_{19}$.
\item $S_{6}$: boundary graph with two trivalent vertices; 18 positive vertices on $\partial S$, 8 four-valent interior vertices, 25 transverse crossings with $\Gamma$, $\chi(S)=8$.
Boundary segments: $t_{1}$ $(9,16)$ $q_{11}$ $(5,12)$ $q_{8}$ $(14,19)$ $q_{7}$ $(6,13)$ $q_{12}$ $(10,15)$ $t_{2}$; $t_{1}$ $(16,3)$ $q_{4}$ $(18,8)$ $q_{5}$ $(15,4)$ $t_{2}$; $t_{1}$ $(3,9)$ $q_{15}$ $(0,1)$ $q_{18}$ $(4,10)$ $t_{2}$.
Node coefficients: $\rho(S_{6})=q_{4}-q_{5}+q_{7}-q_{8}+q_{11}-q_{12}+q_{15}-q_{18}$.
\end{itemize}
\emph{Index.} In the basis above, the node coefficient vectors of $S_1,\dots,S_{6}$ have coordinates $(-1,0,0,0,0,0)$, $(0,0,1,0,0,0)$, $(0,0,0,0,-1,0)$, $(0,0,0,1,0,0)$, $(0,-1,0,0,0,0)$, $(0,0,0,0,0,1)$; the elementary divisors of this matrix are $1,1,1,1,1,1$, so the vectors $\rho(S_i)$ generate $K$.

\subsection{The hexagon examples \texorpdfstring{$X_{88}$}{X88} and \texorpdfstring{$X_{154}$}{X154}}\label{app:data-hex}
For these two examples the rays listed are the vertices of $\Delta$ and the centre of the hexagon; the fan $\Sigma'$ of the
witness of~\ref{hyp:P} may have further rays, which are in the accompanying data. The hexagon is listed in cyclic order
$w_0,\dots,w_5$, and an asterisk marks its node parallelograms. The $\Z$-bases of $K$ are reduced (by the LLL algorithm)
bases of the lattice computed from the circuits.

\emph{The example $X_{88}$.} $\Delta_{88}$ has $21$ vertices, $22$ unit squares and one hexagon. The profile is the $A_1$ subdivision with the parallelograms $\{o,w_1,w_2,w_3\}$ and $\{o,w_4,w_5,w_0\}$; the other two $A_1$ subdivisions also satisfy~\ref{hyp:P}, and their lattices $K$ are identified with this one by matching the pairs of hexagon nodes, up to sign.

\emph{Rays.} $n_{0}=(1,0,0,0)$, $n_{1}=(0,1,0,0)$, $n_{2}=(1,-1,0,0)$, $n_{3}=(0,0,1,0)$, $n_{4}=(0,0,0,1)$, $n_{5}=(0,0,1,-1)$, $n_{6}=(0,0,-1,1)$, $n_{7}=(0,0,0,-1)$, $n_{8}=(0,0,-1,0)$, $n_{9}=(0,-1,0,0)$, $n_{10}=(-1,1,0,0)$, $n_{11}=(-1,1,1,0)$, $n_{12}=(-1,1,0,1)$, $n_{13}=(0,-1,0,-1)$, $n_{14}=(0,-1,-1,0)$, $n_{15}=(-1,0,0,-1)$, $n_{16}=(-1,0,-1,1)$, $n_{17}=(-1,0,-1,0)$, $n_{18}=(-1,0,0,1)$, $n_{19}=(-1,0,1,-1)$, $n_{20}=(-1,0,1,0)$. The centre of the hexagon is $n_{21}=(-1,0,0,0)$. The hexagon is $n_{15}$, $n_{19}$, $n_{20}$, $n_{18}$, $n_{16}$, $n_{17}$.

\emph{Node parallelograms} $q=(a,b,c,d)$: $q_{1}=(11,1,0,3)$, $q_{2}=(12,4,0,1)$, $q_{3}=(12,4,3,11)$, $q_{4}=(13,7,0,2)$, $q_{5}=(14,2,0,8)$, $q_{6}=(14,8,7,13)$, $q_{7}=(15,7,1,10)$, $q_{8}=(15,19,5,7)$, $q_{9}=(15,19,11,10)$, $q_{10}=(17,10,1,8)$, $q_{11}=(17,8,6,16)$, $q_{12}=(17,10,12,16)$, $q_{13}=(17,8,7,15)$, $q_{14}=(20,3,2,9)$, $q_{15}=(9,2,4,18)$, $q_{16}=(18,4,3,20)$, $q_{17}=(19,5,3,20)$, $q_{18}=(19,13,9,20)$, $q_{19}=(16,18,4,6)$, $q_{20}=(16,18,9,14)$, $q_{21}=(18,12,11,20)$, $q_{22}=(14,17,15,13)$, $q_{23}=(21,19,20,18)^{*}$, $q_{24}=(21,16,17,15)^{*}$.

\emph{A $\Z$-basis of $K$} (rank 10): $q_{3}-q_{16}+q_{21}$, $-q_{1}+q_{2}-q_{3}$, $q_{6}-q_{13}-q_{22}$, $q_{4}-q_{5}+q_{6}$, $q_{7}-q_{10}+q_{13}$, $q_{14}+q_{15}+q_{16}$, $q_{2}-q_{5}+q_{10}-q_{12}+q_{15}-q_{20}$, $-q_{1}+q_{4}-q_{7}+q_{9}+q_{14}+q_{18}$, $q_{9}-q_{12}-q_{21}-q_{23}+q_{24}$, $q_{8}-q_{11}+q_{13}-q_{16}+q_{17}-q_{19}$.

\emph{No forced node.} The vector $-q_{1}+q_{2}-2q_{3}-q_{4}+q_{5}+q_{6}-q_{7}+q_{8}-2q_{9}+q_{10}-q_{11}+2q_{12}-2q_{13}-q_{14}-q_{15}-q_{16}+q_{17}-q_{18}-q_{19}+q_{20}+q_{21}-q_{22}+q_{23}-q_{24}$ lies in $K$.

\emph{The example $X_{154}$.} $\Delta_{154}$ has $22$ vertices, $20$ unit squares and one hexagon. The profile is the $A_2$ subdivision with the three parallelograms $\{o,w_0,w_1,w_2\}$, $\{o,w_2,w_3,w_4\}$, $\{o,w_4,w_5,w_0\}$; the other $A_2$ subdivision also satisfies~\ref{hyp:P}.

\emph{Rays.} $n_{0}=(1,0,0,0)$, $n_{1}=(0,1,0,0)$, $n_{2}=(0,0,1,0)$, $n_{3}=(0,0,0,1)$, $n_{4}=(1,0,0,-1)$, $n_{5}=(0,1,-1,0)$, $n_{6}=(0,-1,1,0)$, $n_{7}=(-1,0,0,1)$, $n_{8}=(0,0,0,-1)$, $n_{9}=(0,0,-1,0)$, $n_{10}=(-1,0,0,0)$, $n_{11}=(0,-1,0,0)$, $n_{12}=(-1,0,-1,1)$, $n_{13}=(0,-1,1,-1)$, $n_{14}=(-1,-1,1,0)$, $n_{15}=(-1,-1,0,1)$, $n_{16}=(-1,0,-1,0)$, $n_{17}=(0,-1,0,-1)$, $n_{18}=(-1,-1,1,-1)$, $n_{19}=(-1,-1,-1,1)$, $n_{20}=(-1,-1,0,-1)$, $n_{21}=(-1,-1,-1,0)$. The centre of the hexagon is $n_{22}=(-1,-1,0,0)$. The hexagon is $n_{20}$, $n_{18}$, $n_{14}$, $n_{15}$, $n_{19}$, $n_{21}$.

\emph{Node parallelograms} $q=(a,b,c,d)$: $q_{1}=(13,4,0,6)$, $q_{2}=(19,12,5,9)$, $q_{3}=(19,15,6,11)$, $q_{4}=(17,4,0,11)$, $q_{5}=(17,13,6,11)$, $q_{6}=(17,4,9,21)$, $q_{7}=(17,11,19,21)$, $q_{8}=(18,8,4,13)$, $q_{9}=(18,10,7,14)$, $q_{10}=(12,7,1,5)$, $q_{11}=(16,10,1,5)$, $q_{12}=(16,10,7,12)$, $q_{13}=(16,5,8,20)$, $q_{14}=(20,18,10,16)$, $q_{15}=(18,13,6,14)$, $q_{16}=(19,15,7,12)$, $q_{17}=(20,17,4,8)$, $q_{18}=(21,9,5,16)$, $q_{19}=(20,18,13,17)$, $q_{20}=(21,16,12,19)$, $q_{21}=(22,20,18,14)^{*}$, $q_{22}=(22,14,15,19)^{*}$, $q_{23}=(22,19,21,20)^{*}$.

\emph{A $\Z$-basis of $K$} (rank 9): $q_{2}-q_{18}+q_{20}$, $q_{1}-q_{4}+q_{5}$, $q_{8}-q_{17}+q_{19}$, $q_{10}-q_{11}+q_{12}$, $q_{6}+q_{13}+q_{17}+q_{18}$, $q_{3}-q_{5}-q_{19}-q_{21}+q_{22}$, $-q_{5}+q_{7}+q_{15}-q_{21}+q_{23}$, $q_{9}-q_{12}-q_{20}-q_{21}+q_{23}$, $q_{12}+q_{14}-q_{16}+q_{21}-q_{22}$.

\emph{No forced node.} The vector $-q_{1}+q_{2}+2q_{3}+q_{4}-2q_{5}+2q_{6}-q_{7}+q_{8}-q_{9}-q_{10}+q_{11}+q_{12}+2q_{13}+q_{14}-q_{15}-q_{16}+q_{17}+q_{18}-q_{19}+2q_{20}+q_{21}+q_{22}-2q_{23}$ lies in $K$.

\endgroup
